\documentclass[a4paper,12pt]{article}
\usepackage[utf8]{inputenc}
\usepackage[T1]{fontenc}
\usepackage[french]{babel}
\usepackage{geometry}
\geometry{left=2.5cm, right=2.5cm, top=2.5cm, bottom=2.5cm}

\title{\textbf{Exercice : Gestion de crise Cyber}}
\author{Maël - BTS SIO}
\date{\today}

\begin{document}

\maketitle

\section{Partie 1 -- Analyse de la situation}

\subsection*{Type d'attaque}
Il s'agit d'un \textbf{Rançongiciel} (ou Ransomware). Un pirate a verrouillé les fichiers pour demander de l'argent.

\subsection*{Deux symptômes visibles}
\begin{itemize}
    \item Des fichiers affichent une extension étrange (ex : \texttt{.locked}).
    \item Un message apparaît sur les postes demandant de payer une rançon.
\end{itemize}

\subsection*{Deux risques majeurs pour l'entreprise}
\begin{itemize}
    \item \textbf{Perte d'accès aux fichiers et blocage du travail :} L'entreprise ne peut plus fonctionner car les documents sont illisibles.
    \item \textbf{Perte financière :} Coût de la rançon (si payée) et coût lié à l'arrêt de la production.
\end{itemize}

\section{Partie 2 -- Réactions immédiates}

\subsection*{3 premières actions d'urgence (Confinement)}
\begin{enumerate}
    \item \textbf{Couper immédiatement tous les réseaux :} Débrancher les câbles réseau des postes et des serveurs pour empêcher la propagation.
    \item \textbf{Identifier les postes infectés :} Noter quels postes affichent le message de rançon.
    \item \textbf{Vérifier les sauvegardes :} S'assurer que les copies sont saines et déconnectées du réseau.
\end{enumerate}

\subsection*{Pourquoi ne pas payer la rançon ?}
\begin{itemize}
    \item Rien ne garantit que les pirates fourniront la clé de déchiffrement après le paiement.
    \item Payer finance les criminels et vous désigne comme une cible facile pour une future attaque.
\end{itemize}

\subsection*{Services internes à prévenir en priorité}
\begin{itemize}
    \item La Direction (Chef d'entreprise).
    \item Le service Comptabilité/Finances (pour la gestion du risque financier).
    \item L'ensemble des employés (pour stopper toute utilisation des postes).
\end{itemize}

\section{Partie 3 -- Communication de crise}

\subsection*{Message aux employés}

\begin{center}
    \fbox{
        \begin{minipage}{0.9\textwidth}
            \textbf{OBJET : CONSIGNE URGENTE DE SÉCURITÉ} \\
            
            Nous sommes victimes d'une cyberattaque. Veuillez \textbf{éteindre immédiatement votre ordinateur} (bouton d'alimentation) et ne plus le rallumer. 
            
            N'essayez pas d'accéder à vos fichiers. Le service informatique prend les mesures nécessaires. Nous vous tiendrons informés dès que possible de la reprise du travail. 
            
            Merci de votre coopération.
        \end{minipage}
    }
\end{center}

\subsection*{2 acteurs extérieurs à informer (Légal)}
\begin{itemize}
    \item \textbf{La CNIL :} Obligatoire si des données personnelles (clients ou employés) sont touchées.
    \item \textbf{L'ANSSI ou Cybermalveillance.gouv.fr :} Pour le signalement technique et l'assistance.
\end{itemize}

\section{Partie 4 -- Remédiation et prévention}

\subsection*{Comment vérifier si les sauvegardes sont exploitables ?}
\begin{itemize}
    \item \textbf{Vérifier la date :} S'assurer qu'elles sont assez récentes pour limiter la perte de données.
    \item \textbf{Tester la restauration :} Essayer de restaurer un fichier sur un système isolé (hors réseau) pour vérifier qu'il est lisible et non chiffré.
\end{itemize}

\subsection*{3 mesures pour réduire le risque de récidive}
\begin{itemize}
    \item \textbf{Mots de passe forts et MFA :} Activer la double authentification pour protéger les accès distants et les comptes administrateurs.
    \item \textbf{Sauvegardes déconnectées :} Conserver une copie des données hors ligne (disque dur externe, bande) inaccessible au virus.
    \item \textbf{Sensibilisation des employés :} Former le personnel à détecter le hameçonnage (phishing) et les pièces jointes suspectes.
\end{itemize}

\end{document}